\section{How the Required Value Arrives}
\label{sec:ch11-how-the-value-arrives}
\index{external setter|(}%

Apollo's documentation says that when the router calls a subgraph for a
computed field it \enquote{includes the \code{size} and \code{weight} of the
\code{Product} object passed to the
resolver}~\autocite{apollo2026contributefields}. That is a fair description of
the wire and it is not what happens in the code, and the gap between the two is
this section.

I expected the required field to be a parameter. It is not. Mosaic's reference
resolver has not changed since
chapter~\ref{ch:the-first-cut-extracting-catalog} and still takes one argument
that means anything:

\begin{minted}[fontsize=\footnotesize,breaklines]{csharp}
[ReferenceResolver]
public static Product? ResolveReference(string id, IResolverContext context)
    => ProductKey.TryDecode(id, context, out var productId)
        ? new Product { Id = productId }
        : null;
\end{minted}

That builds a \code{Product} holding a key and nothing else. It never sees the
category. And yet \code{GetShippingCostAsync}, running a moment later, reads
\code{product.Category} and gets \code{KITCHEN}.

\index{HotChocolate!ExternalSetter@\code{ExternalSetter}}%
What sits between them is at the bottom of \code{EntitiesResolver}, read at the
same tag as section~\ref{sec:ch11-what-entities-asks-for}:

\begin{minted}[fontsize=\footnotesize,breaklines]{csharp}
var entity = await resolver(context);

if (entity is not null
    && type.Features.TryGet(out ExternalSetter? externalSetter))
{
    externalSetter.Invoke(context.Schema, type, representation, entity);
}

return entity;
\end{minted}

The reference resolver returns an object, and then something writes onto it out
of the representation. Nothing was passed to anything.
Figure~\ref{fig:ch11-required-field-route} is the whole route, and the middle
of the three boxes inside Mosaic is the step no documentation describes.

\begin{figure}[htbp]
  \centering
  % The route a @requires field's value takes, from the subgraph that owns it to
% the resolver that needs it. Referenced from chapter 11. Bare tikzpicture; the
% figure environment, caption and label live at the call site.
%
% The point of the picture is the last two steps. Everything up to the
% representation is federation; the step after it is HotChocolate writing a
% value onto an object the reference resolver already returned, which is the
% part no documentation describes.
\begin{tikzpicture}[
    font=\small,
    party/.style={draw, rounded corners=2pt, minimum width=24mm,
                  minimum height=8mm, align=center},
    life/.style={gray, dashed},
    msg/.style={-{Stealth[length=2mm]}},
    back/.style={-{Stealth[length=2mm]}, gray},
    inside/.style={draw, rounded corners=2pt, minimum width=40mm,
                   minimum height=7mm, align=center, font=\scriptsize,
                   fill=black!5},
    flow/.style={-{Stealth[length=2mm]}, thick},
    lbl/.style={font=\scriptsize, align=center, fill=white, inner sep=1pt}
  ]

  \node[party] (client) at (0,0)   {Client};
  \node[party] (router) at (3.8,0) {Router};
  \node[party] (cat)    at (8.0,0) {Catalog\\\scriptsize :5101};
  \node[party] (mos)    at (12.2,0) {Mosaic\\\scriptsize :5100};

  \foreach \x in {0, 3.8, 8.0, 12.2} {
    \draw[life] (\x,-0.5) -- (\x,-4.2);
  }

  % 1. the client asks for shippingCost and nothing else
  \draw[msg] (0,-1.0) -- (3.8,-1.0);
  \node[lbl, anchor=south] at (1.9,-0.95) {\code{\{ title shippingCost \}}};

  % 2 and 3. fetch 0, with the required field added
  \draw[msg] (3.8,-2.0) -- (8.0,-2.0);
  \node[lbl, anchor=south] at (5.9,-1.95)
    {\code{title category}\\\code{\_\_typename id}};
  \draw[back] (8.0,-2.9) -- (3.8,-2.9);
  \node[lbl, anchor=south] at (5.9,-2.85) {2 products};

  % 4. the representation
  \draw[msg] (3.8,-3.9) -- (12.2,-3.9);
  \node[lbl, anchor=south] at (8.0,-3.85)
    {\code{\_entities}, each representation \code{\{\_\_typename, id, category\}}};

  % -- what happens inside Mosaic ---------------------------------------------
  %
  % The three boxes sit under the Mosaic lifeline rather than beside it, and
  % they are the only part of the picture that is not a message.
  \node[inside] (ref) at (11.9,-5.4)
    {reference resolver returns\\a \code{Product} holding only the key};
  \node[inside] (set) at (11.9,-6.9)
    {the external setter copies \code{category}\\out of the representation onto it};
  \node[inside] (res) at (11.9,-8.4)
    {\code{shippingCost} reads\\\code{product.Category}};

  \draw[flow] (12.2,-4.2) -- (12.2,-4.7) -- (11.9,-4.7) -- (ref.north);
  \draw[flow] (ref.south) -- (set.north);
  \draw[flow] (set.south) -- (res.north);

\end{tikzpicture}

  \caption{Everything up to the fourth arrow is federation and is written down
  in the specification. The two steps after it are HotChocolate: the reference
  resolver hands back an object that knows only its own key, and a compiled
  setter fills in the rest from the representation before any field resolver
  runs. A property with no setter is where this route quietly ends.}
  \label{fig:ch11-required-field-route}
\end{figure}

\index{HotChocolate!expression trees at schema build}%
That setter is not reflection at request time. \code{ExternalSetterExpression\-Helper}
builds an expression tree once per type when the schema is built and compiles
it, so what runs per entity is a delegate. Reading how it is built is where the
interesting part is.

\subsection*{Key fields are left alone}

\code{BuildFiller} walks the object type's fields and skips any field that
appears in a \code{@key}. The source says why in as many words:

\begin{minted}[fontsize=\footnotesize,breaklines]{text}
// Fields that participate in a '@key' identify the entity and are owned by the reference
// resolver. The representation echoes them back only so the entity can be located, so the
// resolver's returned values are preserved instead of being overwritten.
\end{minted}

Which is the right call. The representation carries the key in whatever form
the other service encoded it, and Mosaic's \code{ProductKey.TryDecode} turned
that into a \code{Guid} on the way in. Copying the encoded string back over the
decoded one would undo the decode.

\subsection*{Everything else is fair game}

\index{external setter!ignores @external on object types}%
Here is the part I did not expect. For an object type, the walk does not check
for \code{@external} at all. It emits a setter call for every field whose
member is a property with a setter, key fields excepted.
\code{TryAddExternalSetter} has two overloads, one for object types and one for
interfaces; the interface one tests each field for an \code{ExternalDirective}
before it emits anything. The object one, which is the one that runs for
\code{Product} and for every other entity in this book, does not.

The class summary states the rule its author had in mind, and it is a
reasonable rule. Its second sentence:

\begin{minted}[fontsize=\footnotesize,breaklines]{text}
The representation is authoritative wherever it carries a value, while resolver
values survive where it is silent.
\end{minted}

\enquote{Wherever it carries a value} is doing more work than it looks like.
The sample makes it concrete. \code{Widget.Name} in the sample's \code{catalog}
subgraph is that subgraph's own field: it carries no \code{@external} there,
and the store is the authority on it. Send three representations, one of them
carrying a \code{name} of its own:

\begin{minted}[fontsize=\footnotesize]{text}
representation 0 -> Hex bolt
representation 1 -> NOT WHAT THE STORE SAYS
representation 2 -> Split washer
\end{minted}

Representation 1 won. Representation 2 carried \code{nonsense: 42}, a field the
type does not have, and was ignored without comment.

\index{security!a subgraph reachable from outside}%
I want to be careful about how large a claim this is. A router never does this.
It sends the key fields and the fields some \code{@requires} named, and nothing
else, so no client of the graph can reach the behaviour. What can reach it is
anything that can post to the subgraph, which in this book's compose file is
anything on the machine. That is a reason to keep subgraphs off the public
network rather than a defect to report, and
chapter~\ref{ch:security-hardening} is where a subgraph stops being reachable
from outside the cluster. It also means that if you ever put a subgraph behind
nothing but a network policy, \code{\_entities} is the field to think hardest
about.

\subsection*{The trap this leaves in your own code}

\index{external@\code{@external}!needs a settable property}%
The practical consequence is small and will cost somebody an afternoon.
\textbf{The property behind an \code{@external} field must have a setter.} Not
\code{init}, which does work here because the compiled expression calls the
setter directly and init-only is a rule the compiler keeps rather than the
runtime, but not get-only. A get-only property compiles, composes, passes every
schema check, and hands your resolver null for every entity in the graph.

What the client then sees is up to the resolver, and the two possibilities are
worth knowing apart. A resolver that treats the missing value as an ordinary
null answers null for every product, quietly, which looks exactly like the
router forgetting to send the field. Mosaic's throws instead, so it answers
\code{Unexpected Execution Error} for every product, which is louder and no
more informative. Both are one property away from working.

Whichever you get, the trace from
section~\ref{sec:ch11-a-field-from-the-other-service} tells the two causes
apart in five seconds. If the representation in \code{extensions.trace} carries
the field and your resolver does not have it, the property is the problem and
the router is not.

\medskip
That is the runtime half. The composer has three opinions about the same two
directives, and one of them changed my code.

\index{external setter|)}%
